\documentclass[a4paper]{article}
\usepackage[paperwidth=19.304cm, paperheight=7.62cm, margin=0.1cm]{geometry}
\usepackage{fontspec,xunicode}
\usepackage{graphicx}
\usepackage{tikz}
\usepackage[absolute]{textpos}
\usepackage{adjustbox}
\setlength{\TPHorizModule}{\paperwidth}\setlength{\TPVertModule}{\paperheight}
\setlength{\parindent}{0pt}

\setmainfont{Linux Libertine}
\setsansfont{Liberation Sans}

% Info
% 
% If you have the above two fonts (and they have the same names on
% your system), then this file should compile in new TeXLives using
% xelatex lanyard.tex
% 
% This makes liberal use of the `textpos` package to place elements on
% a grid (note the difference between parens, brackets, and curly braces):
% \begin{textbock}{width}[opt center point](x position , y position)
%   stuff
% \end{textbock}
% 
% Basic text commands include
% 
% bold: \textbf{text}
% italic: \emph{text}
% size: {\fontsize{pt size}{dist between lines} text}
% font: {\fontspec{font name} text}
% 
% At least on some systems, you can get a list of valid font names by
% running `fc-list` on the command line

\begin{document}

%% Left side

\begin{textblock}{0.2}(0.015,0.03)
  \includegraphics[width=4cm]{fig/main}
\end{textblock}

\begin{textblock}{0.2105}[0,1](0.03,0.90) 
  \begin{small}\textsf{This QR Code \\ Contains My Contact \\ Information.  Scan It.}\end{small}
  % alternately, you can use 
  % {\fontspec{Liberation Sans} This QR Code \\ Contains My Contact \\ Information.  Scan It.}
  % for a one-time use font
\end{textblock}

\begin{textblock}{0.06}(0.12,0.56)
  \includegraphics[width=\textwidth, angle=30]{fig/arrow}
\end{textblock}

\begin{textblock}{0.2105}(0.24,0.03)
  % example of an absolute (36pt) font
  {\fontsize{36}{40} \textbf{Jarrell} } \\
  % example of a relative font specification using one of LaTeX's
  % built-in names: http://texblog.org/2012/08/29/changing-the-font-size-in-latex/
  \begin{Large} Waggoner \end{Large}

  \vspace{2em}

  Student - Graduate
\end{textblock}

\begin{textblock}{0.2105}[0,1](0.24,0.95)
  \includegraphics[width=4cm]{fig/logo}
\end{textblock}

%% Right side

\begin{textblock}{0.1737}[0.5,0.5](0.87,0.4)
  \includegraphics[width=3.3cm]{fig/main}
\end{textblock}

\begin{textblock}{0.1737}[0.5,1](0.87,0.9)
  \begin{center}
    AKYANZBYXM
  \end{center}
\end{textblock}

\begin{textblock}{0.01}[0.5,1](0.75,0.95)
  \begin{center}
    \begin{tikzpicture}
      \draw(0,0) to (0,1);
    \end{tikzpicture}
  \end{center}
\end{textblock}

\begin{textblock}{0.1737}[0.5,0.5](0.63,0.4)
  \includegraphics[width=3.3cm]{fig/main}
\end{textblock}

\begin{textblock}{0.1737}[0.5,1](0.63,0.9)
  \begin{center}
    RKDEILDWKG
  \end{center}
\end{textblock}

\begin{textblock}{0.5}[0.5,0](0.75,0.01)
  \begin{center} for internal use \end{center}
\end{textblock}

% Corner markers

% top left
\begin{textblock}{0.02}[0,0](0.001,0.001)
    \begin{tikzpicture}[x=\textwidth,y=\textwidth]
      \draw(0,0) to (0,1);
      \draw(0,1) to (1,1);
    \end{tikzpicture}
\end{textblock}

% top right
\begin{textblock}{0.02}[1,0](0.999,0.001)
    \begin{tikzpicture}[x=\textwidth,y=\textwidth]
      \draw(0,0) to (0,1);
      \draw(0,1) to (-1,1);
    \end{tikzpicture}
\end{textblock}

% bottom right
\begin{textblock}{0.02}[1,1](0.999,0.999)
    \begin{tikzpicture}[x=\textwidth,y=\textwidth]
      \draw(0,0) to (0,1);
      \draw(0,0) to (-1,0);
    \end{tikzpicture}
\end{textblock}

% bottom left
\begin{textblock}{0.02}[0,1](0.001,0.999)
    \begin{tikzpicture}[x=\textwidth,y=\textwidth]
      \draw(0,0) to (0,1);
      \draw(0,0) to (1,0);
    \end{tikzpicture}
\end{textblock}

% center line
\begin{textblock}{0.01}[0.5,0](0.5,0) 
  \begin{tikzpicture}
    \draw(0,0) to (0,\paperheight);
  \end{tikzpicture}
\end{textblock}

\end{document}
